\documentclass[11pt]{article}
\usepackage{amsmath,amssymb,amsthm,hyperref}
\newtheorem{theorem}{Theorem}
\newtheorem{lemma}{Lemma}
\newtheorem{proposition}{Proposition}
\begin{document}
\title{Projective Cube Lifts and Exact Thin-Shell Supports}
\author{}
\date{November 24, 2026}
\maketitle

\paragraph{Status.}
All results here have complete written proofs. They have not been
formalized in Lean. They repair the geometric support premise of a
Schmitt-style batching argument on the full Boolean cube, but do not
provide flat attention votes and do not prove an improved head bound.
The attention model and its admissibility conditions are unchanged.

\section{One generic positive denominator removes cospherical circuits}

For a finite set $V\subset\mathbb R^m$, write
$\widehat v=(1,v)\in\mathbb R^{m+1}$ and
$Q(v)=\|\widehat v\|^2=1+\|v\|^2$.
An affine circuit is a minimally affinely dependent subset $C$ of $V$.
Fix, for each such $C$, one nonzero dependency
\[
 \sum_{v\in C}\lambda_v\widehat v=0.
\]
Every $\lambda_v$ is nonzero, by minimality, and the dependency space
is one-dimensional. There are finitely many circuits.

\begin{lemma}[Reciprocal linear forms are independent]
Let $u_1,\ldots,u_k$ be pairwise nonproportional nonzero vectors in
$\mathbb R^d$. Then the rational functions
$b\mapsto1/\langle b,u_i\rangle$ are linearly independent over
$\mathbb R$.
\end{lemma}
\begin{proof}
Suppose that their linear combination with coefficients $a_i$ is
zero. Multiplication by the product of the linear forms gives the
polynomial identity
\[
 \sum_i a_i\prod_{j\ne i}\langle b,u_j\rangle=0.
\]
For each $i$, choose $b$ in the hyperplane
$\langle b,u_i\rangle=0$ but in none of the other hyperplanes. Such
a point exists: each other hyperplane intersects this one in a proper
linear subspace, and a real vector space is not a finite union of
proper linear subspaces. Evaluation gives $a_i=0$.
\end{proof}

\begin{theorem}[Generic projective sphere isolation]
Let $V\subset\mathbb R^m$ be finite, and let $\mathcal U$ be a
nonempty open subset of the coefficient vectors $b$ for which
$B_b(v)=\langle b,\widehat v\rangle>0$ for every $v\in V$.
There exists $b\in\mathcal U$ such that the projective points
\[
 y(v)=\frac{\widehat v}{B_b(v)}\in\mathbb R^{m+1}
\]
have the following property: every nonempty affinely independent
subset $S\subset V$ lies on a positive-radius Euclidean sphere
containing no point $y(v)$ with $v\in V\setminus S$.
In particular, for $V=\{0,1\}^m$, $b$ may be chosen with its constant
coefficient and every slope strictly positive.
\end{theorem}

\begin{proof}
For each circuit $C$, the rational function
\[
 F_C(b)=\sum_{v\in C}\frac{\lambda_v Q(v)}{B_b(v)}
\]
is nonzero by the preceding lemma: the vectors $\widehat v$ are
nonproportional and every coefficient $\lambda_vQ(v)$ is nonzero.
After clearing denominators its numerator is a nonzero polynomial.
A nonzero real polynomial vanishes on no nonempty open set. The
product of the finitely many nonzero numerators is nonzero, so some
$b\in\mathcal U$ makes every $F_C(b)$ nonzero.

All $y(v)$ lie in the affine hyperplane $\langle b,y\rangle=1$.
Projective scaling by the nonzero factors $B_b(v)$ preserves the
affine dependence matroid: an original affine dependency with
coefficients $\lambda_v$ becomes the dependency with coefficients
$\mu_v=\lambda_v B_b(v)$, and conversely. In fact
\[
 \sum_v\mu_v=\sum_v\lambda_v B_b(v)=0,
 \qquad \sum_v\mu_v y(v)=\sum_v\lambda_v\widehat v=0.
\]
No transformed circuit is cospherical. If all its points lay on a
sphere with center $c$ and squared radius $r^2$, summing
$\|y(v)\|^2-2\langle c,y(v)\rangle+\|c\|^2-r^2=0$ with weights
$\mu_v$ would give
\[
 0=\sum_v\mu_v\|y(v)\|^2
   =\sum_v\frac{\lambda_vQ(v)}{B_b(v)}=F_C(b),
\]
a contradiction.

Now fix an independent $S$ and one point $v_0\in S$.
Centers of spheres through $y(S)$ are the solutions of the
consistent affine system
\[
 2\langle c,y(v)-y(v_0)\rangle
   =\|y(v)\|^2-\|y(v_0)\|^2\quad(v\in S\setminus\{v_0\}).
\]
Consistency follows from linear independence of the difference
vectors. For $w\notin S$ outside the affine span of $y(S)$,
requiring its inclusion imposes one additional independent affine
equation on this center space, hence excludes a proper affine
hyperplane of centers. For $w$ inside that affine span, $S\cup\{w\}$
contains an affine circuit, and inclusion of $w$ is impossible by
the preceding paragraph. Choose a center outside the finitely many
proper hyperplanes. If $S$ has at least two points, the resulting
radius is automatically positive. If $S$ is a singleton, also avoid
its center point, which is possible in the positive-dimensional
ambient center space. This produces the required sphere.

For the final assertion take $\mathcal U$ to be the strictly
positive coefficient orthant. Every associated $B_b$ is strictly
positive on the cube.
\end{proof}

The theorem does not say that the projective image is in affine
general position. Its affine dependencies are exactly those of the
original cube. It says that no affine circuit is cospherical, which
is precisely the weaker condition needed to isolate an independent
batch by a sphere.

\section{An equivalent rational lift}

\begin{proposition}[Two-ratio exact zero sets]
Let $V\subset\mathbb R^m$ be finite. Fix a function $q:V\to(0,\infty)$
and a positive affine function $D$ on $V$. For generic $B$ in any
nonempty open set of affine functions positive on $V$, every
nonempty affinely independent $S\subset V$ admits an affine $L_S$
such that
\[
 T_S(v)=\frac{q(v)}{B(v)}-\frac{L_S(v)}{D(v)}=0
 \quad\Longleftrightarrow\quad v\in S.
\]
If $q$ is affine and $V$ is the Boolean cube, both terms are actual
admissible affine ratios when $B$ and $D$ have positive constant
coefficients and strictly positive slopes.
\end{proposition}

\begin{proof}
Apply the same reciprocal-independence proof to the finitely many
circuit sums
$\sum_{v\in C}\lambda_v q(v)D(v)/B(v)$, whose coefficients are
nonzero. Put $Z(v)=q(v)D(v)/B(v)$. Since $S$ is affinely independent,
there exist affine functions $L$ interpolating $Z$ on $S$.
For an excluded $w$ outside $\operatorname{aff}(S)$, the additional
equation $L(w)=Z(w)$ cuts a proper affine hyperplane from this
interpolation space. For an excluded $w$ inside the span, that
equation would contradict the nonzero circuit sum. Avoiding the
finitely many proper hyperplanes gives $L_S$. Division by positive
$D(v)$ proves the zero-set statement.
\end{proof}

Thus exact zero sets of arbitrary independent batches cost two
admissible ratios. This is a zero certificate, not an approximate
indicator. Its values away from the batch can have both signs and
arbitrary magnitudes.

\begin{proposition}[Five admissible heads recognize one thin shell]
Let $B$ be an affine function on $\{0,1\}^m$ with strictly positive
constant and slopes, and let $L$ be any affine function. For any
$a>0$ such that
\[
 a^2-\left(\frac1{B(x)}-L(x)\right)^2\ne0
 \qquad(x\in\{0,1\}^m),
\]
the strict signs of this expression are realized by five admissible
affine ratios and a scalar bias.
\end{proposition}
\begin{proof}
Write $L(x)=\ell_0+\sum_i\ell_i x_i$. For $t>0$ set
$D_t(x)=1+\sqrt t\sum_i x_i$ and
\[
 A_t(x)=\frac1t\left(\frac1{B(x)}-\frac1{B(x)+t}\right)
       -\frac{2L(x)}{B(x)}
       +\frac1t\left(\frac{L(x)}{D_t(x)}
                    -\frac{L(x)}{D_t(x)+tL(x)}\right).
\]
This is a sum of five affine ratios. All its denominators have
strictly positive constant and slopes for sufficiently small $t$.
For the only non-immediate denominator, the slopes are
$\sqrt t+t\ell_i>0$ and the constant is $1+t\ell_0>0$; these
inequalities hold simultaneously because there are finitely many
coefficients. Hence every denominator is positive on the entire
cube, not merely on the test set. Algebra gives
\[
 A_t(x)=\frac1{B(x)(B(x)+t)}-\frac{2L(x)}{B(x)}
       +\frac{L(x)^2}{D_t(x)(D_t(x)+tL(x))}.
\]
As $t\downarrow0$, this converges uniformly on the finite cube to
$(1/B-L)^2$. The desired score has a positive minimum absolute
margin, so for sufficiently small $t$ the score $a^2-A_t$ has
exactly its strict signs.
\end{proof}

\begin{proposition}[A batch times arbitrary one-hot columns]
Fix a nonempty independent batch $S\subset\{0,1\}^m$ and a subset
$A\subset\{1,\ldots,q\}$, with $q\ge1$. There exist five admissible
heads on $m+q$ bits and a scalar bias whose score on $(x,e_j)$ is
strictly positive exactly when $(x,j)\in S\times A$.
The denominators are admissible on the full $(m+q)$-cube.
\end{proposition}
\begin{proof}
Apply the rational-lift argument with $q(x)=1$ and $D(x)=1$ to
obtain a positive affine $B_0$ with strictly positive slopes and
an affine $L_0$ such that $T_0=1/B_0-L_0$ vanishes exactly on $S$.
Choose $a>0$ with $4a<|T_0(x)|$ for every $x\notin S$; if there
are no such points, choose any $a>0$.
Choose $0<\eta<B_0(0)$ and put
\[
 B(x,z)=B_0(x)-\eta+\eta\sum_{j=1}^qz_j,
 \qquad L(x,z)=L_0(x)+2a\sum_{j\notin A}z_j.
\]
The denominator $B$ has strictly positive constant and slopes, and
$B(x,e_j)=B_0(x)$. The residual $1/B-L$ equals $0$ on $S\times A$,
equals $-2a$ on $S\times A^c$, and has absolute value greater than
$2a$ at every test point with $x\notin S$. Thus the score
$a^2-(1/B-L)^2$ has the required strict signs on the tests.
The preceding finite-difference construction approximates this
score uniformly on the whole cube. Its test-set margin is positive,
so sufficiently accurate approximation gives the asserted five
heads. No assertion about the signs at non-test inputs is needed.
\end{proof}

\section{A full-cube batching bound}

\begin{lemma}[Vertices in an affine flat]
An affine subspace of dimension $d$ contains at most $2^d$ vertices
of $\{0,1\}^m$.
\end{lemma}
\begin{proof}
On its $d$-dimensional direction space, choose $d$ coordinate
functionals whose restrictions form a basis of the dual space.
Projection onto these coordinates is injective on the affine
subspace. Every Boolean vertex projects into $\{0,1\}^d$.
\end{proof}

\begin{theorem}[Independent monochromatic batch partition]
For $m\ge2$, color the Boolean cube with at most $\kappa$ colors.
Put $k=\lfloor m/2\rfloor+1$. There is a partition into nonempty
monochromatic affinely independent batches, each of size at most
$k$, with total number $K$ satisfying
\[
 K\le\frac{2^m}{k}+\kappa 2^{\lfloor m/2\rfloor}.
\]
For any fixed $a$, if $\kappa\le m^a$ for sufficiently large $m$,
then $K=O(2^m/m)$.
\end{theorem}
\begin{proof}
While a color class has more than $2^{k-2}$ remaining points, its
affine dimension is at least $k-1$, by the flat cardinality lemma,
so it contains $k$ affinely independent points. Remove these as a
batch. There are at most $2^m/k$ such batches. Once no such removal
is possible, each color has at most $2^{k-2}$ points; use singleton
batches for all remaining points. Their total is at most
$\kappa2^{k-2}\le\kappa2^{\lfloor m/2\rfloor}$.
The asymptotic assertion follows because
$m^{a+1}2^{-\lceil m/2\rceil}\to0$.
\end{proof}

Combining this theorem with generic projective sphere isolation
gives a complete geometric repair of the general-position batching
premise for any polynomial number of row colors. Combining it with
the rational-lift proposition gives exact two-ratio zero
certificates for every batch in the same partition.

\section{The remaining vote problem, without changing the input model}

On the Boolean cube, $Q(x)=1+\sum_i x_i$ is affine. A sphere in the
projective image has residual
\[
 \|y(x)-c\|^2-r^2
 =\frac{Q(x)-2B(x)L(x)+\gamma B(x)^2}{B(x)^2},
 \quad L(x)=\langle c,(1,x)\rangle,\quad
 \gamma=\|c\|^2-r^2.
\]
Multiplying by positive $B$ yields
$Q/B-2L+\gamma B$. This exact algebraic simplification explains why
membership signs are inexpensive. It does not flatten those signs.

For the two-ratio version, write
$T_S=q/B-L_S/D$ and
$P_S=qD-L_SB$. A standard narrow bump would require a function of
\[
 T_S^2=\frac{P_S^2}{B^2D^2},
\]
or a reciprocal of $P_S$ when a one-sided construction is possible.
The numerator $P_S$ is generically quadratic, and its reciprocal is
not an affine-ratio head. A simple attempted factor completion
$B(M+\beta)+\alpha(M+\beta)$ can absorb a linear remainder only when
that remainder lies in the span of $1,B,M$; no such condition is
supplied by interpolation of arbitrary batches.

For $K$ batches, flat approximations with absolute error below
$1/(4K)$ would suffice for strict output signs after summation at
threshold $1/2$, as proved in
\texttt{continuation-thin-shell-audit.tex}. Neither genericity nor
the independent-batch partition supplies these approximations.
The common $B,D,q$ across batches could permit a joint cancellation
construction, but no such construction is asserted here.

\begin{proposition}[Raw squared residuals do not supply the joint vote]
Fix common positive affine $B,D$ and affine $q$ on the Boolean
cube. For arbitrary affine $L_s$, real weights $w_s$, real
thresholds $a_s$, and any number of batches, the strict signs of
\[
 F(x)=c+\sum_s w_s\left(a_s^2-
          \left(\frac{q(x)}{B(x)}-\frac{L_s(x)}{D(x)}\right)^2\right)
\]
are the strict signs of a polynomial of degree at most four.
Consequently no such strict score realizes parity on more than
four bits.
\end{proposition}
\begin{proof}
Multiplication by the positive function $B^2D^2$ gives
\[
 B^2D^2F=
 \left(c+\sum_s w_sa_s^2\right)B^2D^2
 -\sum_s w_s(qD-L_sB)^2.
\]
Every term has degree at most four. To recall the parity obstruction
directly, write parity as $\chi(x)=(-1)^{\sum_i x_i}$. The uniform
average of $\chi$ times every monomial of degree less than $m$ is
zero, by summing over an absent coordinate. Hence the average of
$\chi P$ is zero for any degree-less-than-$m$ polynomial $P$.
Strict sign representation would instead make every summand
$\chi(x)P(x)$ positive, a contradiction.
\end{proof}

This last proposition concerns the displayed strict limiting score.
It does not exclude a finite-difference construction in which all
lower-order terms cancel identically and the signs are carried by a
higher-order term. Such a mechanism would need its own construction
and uniform strict-sign proof.

\end{document}
